% ============================================================
%  Probabilidade — Apostila Premium | PratiqueJá
%  Engine: XeLaTeX
% ============================================================
\documentclass[12pt]{article}
\usepackage{../../pratiqueja}

\newcommand{\hlm}[1]{{\color{darkblue}#1}}

\setsubject{Probabilidade}

\begin{document}
\pagenumbering{arabic}
\setstretch{1.2}
\color{bodytext}

\heroheader{Probabilidade}%
           {Espaço amostral, regras e distribuição binomial}%
           {Matemática · Avançado}

\setlength{\columnsep}{18pt}
\setlength{\columnseprule}{0.5pt}
\def\columnseprulecolor{\color{exborder}}

\begin{multicols}{2}

\TW{Def}{darkblue}{Definição}{O que é probabilidade?}

\regra{
Probabilidade é a medida numérica da chance de um evento ocorrer, variando de $0$ (evento impossível) a $1$ (evento certo).
}

\nota{\textbf{Representação:} fração, decimal ou porcentagem. Ex.: $\dfrac{1}{4} = 0{,}25 = 25\%$.}

\TW{1}{darkblue}{Espaço Amostral}{Resultados, eventos e nomenclatura}

\regra{
O \hl{espaço amostral} $\Omega$ é o conjunto de \emph{todos} os resultados possíveis. Um \hl{evento} $A$ é qualquer subconjunto de $\Omega$.
}

\exemp{
Dado: $\Omega = \{1,2,3,4,5,6\}$, $n(\Omega) = 6$

Evento "par": $A = \{2,4,6\}$, $n(A) = 3$

Evento "maior que 4": $B = \{5,6\}$, $n(B) = 2$

\textbf{Evento certo:} $\Omega$ — $P(\Omega) = 1$

\textbf{Evento impossível:} $\varnothing$ — $P(\varnothing) = 0$

\textbf{Eventos exclusivos:} $A \cap B = \varnothing$
}

\TW{2}{darkblue}{Cálculo de Probabilidade}{Fórmula clássica de Laplace}

\regra{
Válida quando todos os resultados são \hl{igualmente prováveis}:
\[
  P(A) = \dfrac{n(A)}{n(\Omega)}
\]
}

\exemp{
\textbf{P(face par) num dado:}

$n(A) = 3$ \quad $n(\Omega) = 6$

$\hlm{P(A) = \dfrac{3}{6} = \dfrac{1}{2} = 50\%}$
}

\TW{$A^c$}{badgepurple}{Eventos Complementares}{Probabilidade do evento não ocorrer}

\regra{
\[
  P(A^c) = 1 - P(A)
\]
$A^c$ = pontos de $\Omega$ que \textbf{não} pertencem a $A$.
}

\exemp{
Pote: 12 pretos, 16 azuis, 2 vermelhos ($n(\Omega) = 30$).

$P(\text{vermelho}) = \dfrac{2}{30} = \dfrac{1}{15}$

$\hlm{P(\text{não vermelho}) = 1 - \dfrac{1}{15} = \dfrac{14}{15} \approx 93{,}3\%}$
}

\nota{\textbf{Estratégia:} quando $P(A)$ tem muitos casos favoráveis, calcule $P(A^c)$ e subtraia de 1.}

\TW{$\geq$1}{badgepurple}{Probabilidade de ``Ao Menos Um''}{Aplicação direta do evento complementar}

\regra{
\[
  P(\text{ao menos um}) = 1 - P(\text{nenhum})
\]
}

\exemp{
\textbf{Moeda lançada 3 vezes — ao menos uma cara:}

$P(\text{nenhuma cara}) = \left(\dfrac{1}{2}\right)^3 = \dfrac{1}{8}$

$\hlm{P(\text{ao menos uma cara}) = \dfrac{7}{8} = 87{,}5\%}$

\textbf{Dado lançado 4 vezes — ao menos um 6:}

$P(\text{não sair 6}) = \dfrac{5}{6}$ por lançamento

$P(\text{nenhum}) = \left(\dfrac{5}{6}\right)^4 = \dfrac{625}{1296}$

$\hlm{P(\text{ao menos um 6}) = 1 - \dfrac{625}{1296} = \dfrac{671}{1296} \approx 51{,}8\%}$
}

\nota{\textbf{Atenção:} funciona para lançamentos/tentativas \emph{independentes}. Para eventos dependentes, use combinações.}

\TW{P|A}{darkblue}{Probabilidade Condicional}{Probabilidade de B dado que A ocorreu}

\regra{
\[
  P(B \mid A) = \dfrac{P(A \cap B)}{P(A)}, \quad P(A) \neq 0
\]
}

\exemp{
Urna: 3 bolas vermelhas + 2 azuis.

Sabendo que a 1ª foi vermelha, restam 4 bolas (2V, 2A):

$\hlm{P(B \mid A) = \dfrac{2}{4} = \dfrac{1}{2} = 50\%}$
}

\TW{$\cup$}{darkblue}{Regra da Soma}{Probabilidade de A ou B ocorrerem}

\regra{
\textbf{Eventos mutuamente exclusivos} ($A \cap B = \varnothing$):
\[
  P(A \cup B) = P(A) + P(B)
\]

\textbf{Eventos não exclusivos:}
\[
  P(A \cup B) = P(A) + P(B) - P(A \cap B)
\]
}

\exemp{
\textbf{P(sair 2 ou 5) no dado:}

Exclusivos: $\hlm{P = \dfrac{1}{6} + \dfrac{1}{6} = \dfrac{1}{3}}$

\textbf{Turma: 18 estudam Mat. $(A)$, 12 Fís. $(B)$, 5 ambas:}

$P(A \cup B) = \dfrac{18}{30} + \dfrac{12}{30} - \dfrac{5}{30} = \hlm{\dfrac{25}{30} = \dfrac{5}{6}}$
}

\TW{$\cap$}{darkblue}{Regra do Produto}{Probabilidade de A e B ocorrerem juntos}

\regra{
\textbf{Eventos independentes:}
\[
  P(A \cap B) = P(A) \cdot P(B)
\]
\textbf{Eventos dependentes:}
\[
  P(A \cap B) = P(A) \cdot P(B \mid A)
\]
}

\exemp{
\textbf{Cara na moeda E 4 no dado (independentes):}

$P = \dfrac{1}{2} \cdot \dfrac{1}{6} = \hlm{\dfrac{1}{12}}$

\textbf{Caixa: 5 vermelhas + 3 azuis — vermelho depois azul (s/ reposição):}

$P(A) = \dfrac{5}{8}$ \quad $P(B \mid A) = \dfrac{3}{7}$

$\hlm{P(A \cap B) = \dfrac{5}{8} \cdot \dfrac{3}{7} = \dfrac{15}{56}}$
}

\nota{\textbf{Sem reposição} → dependentes: use $P(A) \cdot P(B|A)$.

\textbf{Com reposição} → independentes: use $P(A) \cdot P(B)$.}

\TW{$C_n$}{badgepurple}{Probabilidade e Combinatória}{Usando combinações para contar casos}

\regra{
Em retiradas simultâneas, use combinações para contar:
\[
  P(A) = \dfrac{C_{n,k}}{C_{N,k}}, \quad C_{n,k} = \dfrac{n!}{k!\,(n-k)!}
\]
}

\exemp{
Urna: 4 brancas + 3 pretas. Retira-se 2 simultaneamente.

$n(\Omega) = C_{7,2} = 21$ \quad $n(A) = C_{4,2} = 6$

$\hlm{P(\text{2 brancas}) = \dfrac{6}{21} = \dfrac{2}{7} \approx 28{,}6\%}$

\textbf{P(cores diferentes):}

$n(B) = C_{4,1} \cdot C_{3,1} = 12$

$\hlm{P(B) = \dfrac{12}{21} = \dfrac{4}{7} \approx 57{,}1\%}$
}

\TW{B(n,p)}{badgepurple}{Distribuição Binomial}{Experimentos de Bernoulli repetidos}

\regra{
Quando um experimento com $P(\text{sucesso}) = p$ é repetido $n$ vezes de forma independente, a probabilidade de exatamente $k$ sucessos é:
\[
  P(X = k) = \binom{n}{k} p^k (1-p)^{n-k}
\]
}

\exemp{
\textbf{Dado lançado 5 vezes — exatamente 2 faces "6":}

$n=5$, $k=2$, $p=\frac{1}{6}$

$P = C_{5,2} \cdot \left(\frac{1}{6}\right)^2 \cdot \left(\frac{5}{6}\right)^3 = 10 \cdot \frac{125}{7776}$

$\hlm{P \approx 16{,}1\%}$

\textbf{Fábrica: 10\% defeituosas, lote de 4. P(nenhuma defeituosa):}

$P(X=0) = (0{,}90)^4 = \hlm{0{,}6561 = 65{,}61\%}$
}

\nota{\textbf{Quando usar:} experimentos (1) independentes, (2) com mesmo $p$ de sucesso e (3) $n$ fixo.}

\TW{!}{badgepurple}{Dicas e Erros Comuns}{O que mais cai em vestibular e concurso}

\exemp{
\textbf{Com vs. sem reposição:}

Sem reposição = dependentes → use $P(A) \cdot P(B|A)$

Com reposição = independentes → use $P(A) \cdot P(B)$

\textbf{``Ao menos um'' → use complementar:}

$P(\text{ao menos um}) = 1 - P(\text{nenhum})$

\textbf{Regra da Soma — não esqueça a interseção:}

$P(A \cup B) = P(A) + P(B) - P(A \cap B)$

\textbf{Independência $\neq$ Exclusividade:}

Exclusivos com $P > 0$ são \emph{dependentes}.

Independência: $P(A \cap B) = P(A) \cdot P(B)$.
}

\end{multicols}

\end{document}
